\scanpage{37}
Lemme de Poincaré vrai pour les $k$-algèbres suivantes ($k$ corps de car. nulle).
\[
\left\{\begin{aligned}
&k[[x]][\unsure{Y|X}]/(Y^p-x^q)\\
&k[[x]][Y]/(Y^2+P(x)Y+Q(x))\\
&k[x_1,\ldots x_n]/(x_1^{p_1},x_2^{p_2},\ldots,x_n^{p_n})\qquad\text{p. ex. }k[x]/(x^p)\\
&k[x,y]/\mathfrak m^2\\
&k[x,y]/(x^3,y^2,x^2y)\\
&k[x,y]/(x^3,y^2,xy)
\end{aligned}\right.
\]
Si $A$ \deleted{\unsure{Artinien}} local sur $k$, ext. résiduelle triviale, alors pour tt $x\in\mathfrak m$ avec $d_{A/k}x=0$, on a $x\in\mathfrak m^3$.

\figure{\begin{tikzpicture}
\node[anchor=west] at (0,0) {Ann. Éc. N. Sc.};
\node[anchor=west] at (4.3,0.13) {$65\ \unsure{ou|à}\ 66$};
\node[anchor=west] at (0.4,-0.65) {ou};
\node[anchor=west] at (1.6,-0.65) {J. de Liouville};
\draw (4.15,0.4)--(4.1,-0.14)--(5.9,-0.14);
\draw (5.85,0.4)--(6.02,0.4)--(5.96,-0.75)--(5.79,-0.75);
\draw (5.89,0.38)--(5.84,-0.73);
\end{tikzpicture}}
\mspage{36}
